(1) 方程为 $-\hbar^2\psi''/(2m)-g\delta(x)\psi=E\psi$。对 $[-\epsilon,\epsilon]$ 积分并令 $\epsilon\to0$，有限能量项积分消失而点势保留，得到 $\psi'(0^+)-\psi'(0^-)=-2mg\,\psi(0)/\hbar^2$。波函数本身连续，否则二阶导会包含无法由单一 $\delta$ 势抵消的更高阶分布。
(2) 负能束缚解在两侧必须指数衰减，且吸引性点势只有偶宇称态，令 $\psi_b(x)=A e^{-\kappa|x|}$。跃变条件给 $-2\kappa A=-2mgA/\hbar^2$，即 $\kappa=mg/\hbar^2$；所以 $E_b=-\hbar^2\kappa^2/(2m)=-mg^2/(2\hbar^2)$。由 $\int_{-\infty}^{\infty}e^{-2\kappa|x|}dx=1/\kappa$ 得 $A=\sqrt\kappa$，不存在第二个可归一化独立负能解。
(3) 取 $x<0$ 的 $\psi=e^{ikx}+r e^{-ikx}$、$x>0$ 的 $\psi=t e^{ikx}$。连续性给 $t=1+r$，跃变条件给 $2ik r=-2\alpha t$，$\alpha=mg/\hbar^2$。解得 $r=-\alpha/(ik+\alpha)$、$t=ik/(ik+\alpha)$。这里两侧粒子色散相同，概率流率与振幅平方直接成正比，无需额外速度比因子。
(4) 概率系数为 $R=|r|^2=\alpha^2/(k^2+\alpha^2)$、$T=|t|^2=k^2/(k^2+\alpha^2)$，可直接验证 $R+T=1$。$g\to0$ 时 $\alpha\to0$ 得 $T\to1$，强吸引 $g\to\infty$ 时固定正能粒子仍趋于全反射。解析延拓 $t(k)=ik/(ik+\alpha)$ 在 $k=i\alpha$ 处出现上半复波数平面的极点，正对应 $E_b=-\hbar^2\alpha^2/(2m)$。
